\documentclass[aspectratio=169,10pt,english]{beamer}

\input{../../../_preambulo_beamer.tex}
\input{../../../_comandos_beamer.tex}

\selectlanguage{english}

\title{MA1001B - Statistical Modeling for Decision Making}
\subtitle{Lesson 42: Chi-Square Test of Homogeneity}
\author{}
\institute{Tecnol\'ogico de Monterrey}
\date{}

\begin{document}

\begin{frame}[plain]
  \vspace{1.2cm}
  {\Large\bfseries MA1001B - Statistical Modeling for Decision Making\par}
  \vspace{0.7cm}
  {\large Lesson 42: Chi-Square Test of Homogeneity\par}
  \vspace{0.7cm}
  {\color{TecRojo}\rule{\textwidth}{1.2pt}}
  \vspace{0.8cm}

  MA1001B Faculty\\[0.3cm]
  Term\\[0.3cm]
  Tecnol\'ogico de Monterrey
\end{frame}

\begin{frame}{The Shared-Mix Question}
  \begin{alertblock}{Guiding question}
    Do three independently sampled warehouses share one ticket-channel mix?
  \end{alertblock}

  \vspace{0.6em}
  By the end of this lesson, you can:
  \begin{itemize}
    \item treat row totals as fixed by the sampling plan;
    \item compute a homogeneity $\chi^2$ with $df=(r-1)(c-1)$;
    \item decide whether the sites share one mix;
    \item contrast that story with a test of independence.
  \end{itemize}

  \vfill
  \scriptsize All ticket counts used today are synthetic. No real company data are used.
\end{frame}

\begin{frame}{Three Planned Samples, Same Categories}
  \small
  \begin{center}
    \begin{tabular}{lrrrr}
      \toprule
      \textbf{Site} & \textbf{Phone} & \textbf{Chat} & \textbf{Email} & \textbf{$n$} \\
      \midrule
      North & 60 & 40 & 20 & 120 \\
      Central & 35 & 40 & 25 & 100 \\
      South & 25 & 30 & 35 & 90 \\
      \midrule
      Total & 120 & 110 & 80 & 310 \\
      \bottomrule
    \end{tabular}
  \end{center}

  \vspace{0.3em}
  \begin{exampleblock}{Homogeneity hypotheses}
    $H_0$: the three sites share one mix versus $H_1$: at least one site
    differs, with $\alpha=0.05$.
  \end{exampleblock}
\end{frame}

\begin{frame}[fragile]{Step 1: Three Independent Mixes}
  \begin{columns}[T]
    \begin{column}{0.42\textwidth}
      \small
      \begin{lstlisting}
shares = counts / n_site
      \end{lstlisting}

      \begin{block}{Verified output}
        North phone share $=0.500000$\\
        Central phone share $=0.350000$\\
        South phone share $=0.277778$
      \end{block}
    \end{column}
    \begin{column}{0.58\textwidth}
      \centering
      \includegraphics[width=\textwidth]{../figures/homogeneity_01_three_samples.png}
      \vspace{0.2em}
      \scriptsize Within-site mixes from Step 1
    \end{column}
  \end{columns}
\end{frame}

\begin{frame}{Homogeneity Uses the Same $\chi^2$ Arithmetic}
  \small
  Expected counts still use the product of margins:
  \[
    E_{ij}=\frac{(\text{row total}_i)(\text{column total}_j)}{n}.
  \]
  \[
    \chi^2=17.812605,\qquad df=(3-1)(3-1)=4,\qquad p=0.001343.
  \]
  Because $\chi^2>\chi^{2*}=9.487729$ and $p<\alpha$, the test rejects the
  claim that the three sites share one mix.
\end{frame}

\begin{frame}[fragile]{Step 2: Homogeneity Decision}
  \begin{columns}[T]
    \begin{column}{0.40\textwidth}
      \small
      \begin{lstlisting}
chi2, p, df, E = chi2_contingency(table)
      \end{lstlisting}

      \begin{block}{Verified output}
        $\chi^2=17.812605$, $df=4$\\
        $p=0.001343$\\
        Decision: reject $H_0$
      \end{block}
    \end{column}
    \begin{column}{0.60\textwidth}
      \centering
      \includegraphics[width=\textwidth]{../figures/homogeneity_02_chi_square.png}
      \vspace{0.2em}
      \scriptsize Homogeneity cell contributions from Step 2
    \end{column}
  \end{columns}
\end{frame}

\begin{frame}{Step 3: Same Statistic, Different Sampling Story}
  \begin{columns}[T]
    \begin{column}{0.42\textwidth}
      \small
      Homogeneity: three planned samples. Row totals $120$, $100$, and $90$
      were chosen first.

      \vspace{0.4em}
      Independence: one sample of $310$ tickets later classified by site and
      channel.

      \vspace{0.5em}
      \textbf{Limit:} both stories produce $\chi^2=17.812605$ and
      $p=0.001343$. The arithmetic is shared; the $H_0$ is not.
    \end{column}
    \begin{column}{0.58\textwidth}
      \centering
      \includegraphics[width=\textwidth]{../figures/homogeneity_03_sampling_story_limit.png}
      \vspace{0.2em}
      \scriptsize Two sampling stories from Step 3
    \end{column}
  \end{columns}
\end{frame}

\begin{frame}{Concept Check}
  \begin{enumerate}
    \item Why are the row totals $120$, $100$, and $90$ part of the design?
    \item Compute $df$ for three sites and three categories.
    \item If $p=0.001343$ and $\alpha=0.05$, what is the homogeneity
      decision?
    \item How can the same table be an independence problem instead?
  \end{enumerate}

  \vspace{0.6em}
  \begin{block}{Connection to Lesson 43}
    The session can continue directly into experimental design: factors,
    randomization, and why an observational site is not a treatment.
  \end{block}

  \vfill
  \scriptsize Source: OpenStax, \textit{Introductory Business Statistics 2e},
  Chapter 11, Section 11.5. CC BY-NC-SA.
\end{frame}

\appendix



\end{document}
